\section{PX4 Software-in-the-Loop Flights in Gazebo}

\begin{figure*}[t]\centering
\begin{minipage}[b]{0.14\textwidth}\centering\includegraphics[width=\linewidth]{gazebo_a.jpg}\\[2pt]{\small (a) Thick fog}\end{minipage}\hspace{2mm}%
\begin{minipage}[b]{0.245\textwidth}\centering\includegraphics[width=\linewidth]{gazebo_b.jpg}\\[2pt]{\small (b) Voxels and planned path}\end{minipage}\hspace{2mm}%
\begin{minipage}[b]{0.245\textwidth}\centering\includegraphics[width=\linewidth]{gazebo_c.jpg}\\[2pt]{\small (c) Magpie near the drone}\end{minipage}\hspace{2mm}%
\begin{minipage}[b]{0.245\textwidth}\centering\includegraphics[width=\linewidth]{gazebo_d.jpg}\\[2pt]{\small (d) Pad at touchdown}\end{minipage}
\caption{Frames from the PX4 flights in Gazebo.}\label{fig:gazebo}\end{figure*}

\textbf{Weather.} With GNSS, the rule deciding and the marker guiding touchdown, we flew ten flights in each of five skies on the same seeds (Figure~\ref{fig:gazebo}(a)). Wind of 7~m/s with gusts to about 10~m/s cost no landing. Mist and thick fog cost one and three landings. Every failure was a safe abort: the drone lost its marker in a fog patch and searched away, once only 0.68~m from its own pad. With commit-on-lock, mist and fog both landed 10 of 10, and the median miss was 0.33 to 0.37~m in every sky. The marker was first read at about 51~m in every sky.

\textbf{Unmapped crane.} We placed a crane at ten positions on the approach and flew the rig with GNSS. A planner that knew only the survey crashed into the crane in 8 of 10 flights. A planner that looked ahead on the voxel map passed all ten but came within 0.99~m. The rule with the local 3D planner passed all ten, never closer than 2.59~m (median 3.42~m).

\textbf{Magpie with the live detector.} With GNSS, the drone knew the magpie only through the up-camera rig's live detector and tracker. The planner and a learner trained with that rig's model each landed 5 of 9 flights. These flights span three revisions of the landing code, so they show only that detection works in the loop.

\begin{table*}[t]
\centering\small
\setlength{\tabcolsep}{5pt}
\caption{PX4 flights in Gazebo with GNSS denied, the last flight of each scenario. Errors against Gazebo's truth, median / 90th percentile. Miss: touchdown distance from the pad center. Flying learner: wind seed 0. Learner: trained with the up-camera rig's model, not Table~\ref{tab:ladder}'s, flying the probability-weighted mean of its commands.}
\label{tab:gnssfree}
\begin{tabular}{llcccc}
\toprule
Scenario & Decider & GNSS off (s) & Visual fix error (m) & PX4 error (m) & Miss (m) \\
\midrule
Calm & Rule & 110 & 2.13 / 4.95 & 2.32 / 3.89 & 0.06 \\
Wind, W10 9~m/s & Flying learner & 117 & 2.12 / 4.71 & 1.71 / 6.68 & 0.19 \\
Crane on the route & Rule, local 3D planner & 117 & 2.28 / 5.12 & 1.85 / 3.89 & 0.04 \\
Magpie & Learner & 116 & 2.21 / 5.12 & 1.67 / 3.55 & 0.04 \\
Magpie, up camera added & Learner & 116 & 2.17 / 5.11 & 1.44 / 3.52 & 0.08 \\
Mission to another pad & Rule & 107 & 2.26 / 5.17 & 3.13 / 4.82 & 0.05 \\
\bottomrule
\end{tabular}
\end{table*}


\textbf{GNSS-denied deliveries.} Table~\ref{tab:gnssfree} lists six flights with GNSS denied over the suburb, the last flight of each scenario. The drone cruised at 75~m on a planned route, and we denied GNSS once PX4 had accepted four visual fixes within 15~s. The visual localization module matched each downward image against satellite imagery captured in 2021. The simulated ground and roofs carried a 2023 capture of the same area, so the map never equals what the camera sees. Per flight, the median fix error was 2.1 to 2.3~m, the 90th percentile 4.7 to 5.2~m, and 97.8 to 98.5\% of fixes fell within 10~m. PX4's fused position had a median error of 1.4 to 3.1~m. On the laptop that ran the simulation, a visual fix was ready a median of 11 to 41~ms after its image was taken. Once the pad's markers were readable they took over, and every flight landed within 0.19~m of the pad center. The touchdown accuracy therefore comes from the markers, and the visual fixes had to bring the drone close enough to read them.

These six flights are integration demonstrations, not a landing rate. Before them, and after short pipeline tests, we flew 18 development flights in the same scenarios. Six failed: a hard landing, a crash in wind, lost visual fixes, a crash into the crane, a timeout and a landing 659~m off the pad. Each led to a fix: a downward rangefinder, marker fixes kept until touchdown with an outlier check, and one simulator per computer. Weighted A* fixed the crane crash by cutting a plan from about 2~s to 17~ms.

We then repeated the calm flight ten times on one frozen code version, changing only the landing heading. All ten landed, 0.04 to 0.07~m from the pad center, with GNSS denied for 109 to 130~s and PX4's median error 2.0 to 2.9~m. As with the weather flights, 10 of 10 shows only that the landing rate exceeds 69\% at 95\% confidence, and calm air is the easiest case.

The rig's cameras and the detector shared one laptop GPU (RTX 5070), and we slowed each magpie landing to a tenth of real time for the detector. Timed separately on that GPU, the detector took a median of 54 to 65~ms per set of four or five pictures. The slowest set took 157~ms, inside the 0.25~s reflex period. It ran beside the slowed simulation without pausing it, so its delay cost the drone a tenth of the simulated time it would in real time. These flights validate the software integration, not real-time operation on a drone's own computer.
